\chapter{Registro de decisiones arquitectónicas}
\label{anexo:decisiones-arquitectonicas}

Este anexo documenta las principales decisiones de arquitectura adoptadas durante el desarrollo del prototipo mediante Architecture Decision Records (ADR). El objetivo es conservar, además de la decisión final, el contexto que la motivó, las alternativas consideradas y sus principales consecuencias. Los ADR describen decisiones técnicas del prototipo y pueden revisarse si cambian sus requerimientos o condiciones de despliegue.

\section*{ADR-001 --- Separación de Frontend, Backend y AI Module}
\phantomsection
\addcontentsline{toc}{section}{ADR-001 --- Separación de Frontend, Backend y AI Module}
\label{adr:001-separacion-componentes}

\textbf{Estado:} aceptado.

\textbf{Contexto.} La solución combina una interfaz interactiva, funciones de autenticación y persistencia, y un runtime de procesamiento de imágenes e inteligencia artificial. Estos componentes presentan responsabilidades, dependencias y ciclos de evolución diferentes.

\textbf{Decisión.} Mantener Frontend, Backend y AI Module como componentes diferenciados y alojados en repositorios independientes. Cada componente conserva sus propias dependencias, configuración y pruebas, y la integración se realiza mediante contratos explícitos.

\textbf{Alternativas consideradas.} Unificar el sistema en un único servicio o repositorio; integrar el runtime de IA directamente en el Backend.

\textbf{Consecuencias.} La separación permite modificar modelos o componentes de interfaz sin incorporar sus dependencias al resto del sistema y facilita pruebas y despliegues independientes. Como contrapartida, exige mantener contratos compatibles y coordinar cambios entre repositorios.

\section*{ADR-002 --- Spring Boot para el Backend y Python/FastAPI para el AI Module}
\phantomsection
\addcontentsline{toc}{section}{ADR-002 --- Spring Boot para el Backend y Python/FastAPI para el AI Module}
\label{adr:002-spring-boot-fastapi}

\textbf{Estado:} aceptado.

\textbf{Contexto.} El Backend concentra responsabilidades de producto —autenticación, autorización, persistencia, auditoría, reglas de negocio y exposición de contratos— mientras que el módulo de IA depende del ecosistema utilizado para procesamiento de imágenes, entrenamiento e inferencia.

\textbf{Decisión.} Utilizar Spring Boot y Java para el Backend de producto, y Python con FastAPI para el AI Module.

\textbf{Alternativas consideradas.} Implementar ambos servicios con Python y FastAPI; ejecutar también la inferencia desde Java.

\textbf{Consecuencias.} Cada componente utiliza un stack alineado con sus responsabilidades y dependencias. PyTorch y SimpleITK permanecen aislados en el servicio de IA, mientras las funciones de producto se mantienen independientes del runtime científico. La principal contrapartida es mantener dos stacks tecnológicos y su integración HTTP.

\section*{ADR-003 --- Backend como única frontera pública hacia el AI Module}
\phantomsection
\addcontentsline{toc}{section}{ADR-003 --- Backend como única frontera pública hacia el AI Module}
\label{adr:003-backend-frontera-publica-ai}

\textbf{Estado:} aceptado.

\textbf{Contexto.} Los resultados de inferencia contienen información técnica y referencias internas que no necesariamente deben exponerse al cliente web. Además, las operaciones sobre estudios y resultados requieren autenticación, autorización, persistencia y trazabilidad.

\textbf{Decisión.} El Frontend se comunica con el Backend y no consume directamente el AI Module. El Backend autentica y autoriza las operaciones, orquesta las llamadas al servicio de IA, persiste los resultados y expone al Frontend únicamente el contrato público correspondiente.

\textbf{Alternativas consideradas.} Permitir que el Frontend consuma directamente los endpoints del AI Module.

\textbf{Consecuencias.} Se mantiene una única frontera de seguridad y un contrato público controlado, y el runtime de IA puede evolucionar sin convertirse en una dependencia directa de la interfaz. Como contrapartida, el Backend incorpora responsabilidades adicionales de orquestación y existe un salto de comunicación adicional.

\section*{ADR-004 --- PostgreSQL para persistencia transaccional}
\phantomsection
\addcontentsline{toc}{section}{ADR-004 --- PostgreSQL para persistencia transaccional}
\label{adr:004-postgresql-persistencia}

\textbf{Estado:} aceptado.

\textbf{Contexto.} El producto debe conservar relaciones entre usuarios, roles, sujetos de-identificados, estudios, corridas, modelos, mediciones, revisiones, assets y eventos de auditoría.

\textbf{Decisión.} Utilizar PostgreSQL como base de datos transaccional del sistema.

\textbf{Alternativas consideradas.} Persistencia basada únicamente en archivos, almacenamiento en memoria o una base de datos no relacional como mecanismo principal.

\textbf{Consecuencias.} El modelo relacional permite aplicar restricciones de integridad y mantener relaciones consistentes entre las entidades del flujo. Como contrapartida, requiere gestionar el esquema, sus migraciones y la disponibilidad de la instancia de base de datos.

\section*{ADR-005 --- PyTorch como runtime de inferencia y ONNX Runtime diferido}
\phantomsection
\addcontentsline{toc}{section}{ADR-005 --- PyTorch como runtime de inferencia y ONNX Runtime diferido}
\label{adr:005-pytorch-onnx-runtime}

\textbf{Estado:} aceptado para PyTorch; ONNX Runtime diferido.

\textbf{Contexto.} Los modelos utilizados por el prototipo fueron entrenados, evaluados y registrados mediante PyTorch, y los checkpoints actuales dependen de dicho runtime.

\textbf{Decisión.} Mantener PyTorch como runtime de inferencia durante esta etapa. Una eventual migración a ONNX Runtime deberá realizarse únicamente después de verificar la equivalencia de las salidas de los modelos convertidos.

\textbf{Alternativas consideradas.} Migrar inmediatamente los modelos a ONNX Runtime o reconstruir la inferencia utilizando otro ecosistema.

\textbf{Consecuencias.} Se evita introducir una conversión adicional sobre modelos ya evaluados y se conserva la reproducibilidad entre entrenamiento e inferencia. Como contrapartida, el despliegue mantiene las dependencias y requerimientos de recursos propios de PyTorch.

\section*{ADR-006 --- Separación entre código fuente y artifacts de datos y modelos}
\phantomsection
\addcontentsline{toc}{section}{ADR-006 --- Separación entre código fuente y artifacts de datos y modelos}
\label{adr:006-separacion-codigo-artifacts}

\textbf{Estado:} aceptado.

\textbf{Contexto.} Los checkpoints, datasets, volúmenes y demás artifacts poseen tamaños y ciclos de vida distintos de los del código fuente y requieren controles específicos de acceso, integridad y versionado.

\textbf{Decisión.} Mantener el código y la configuración versionable en los repositorios del proyecto, y almacenar modelos, datasets y artifacts de mayor tamaño en almacenamiento externo controlado. Los modelos se identifican mediante manifests, versiones y hashes verificables.

\textbf{Alternativas consideradas.} Incorporar directamente los archivos de modelos y datasets dentro de los repositorios o depender exclusivamente del filesystem local de cada servicio.

\textbf{Consecuencias.} Los repositorios permanecen manejables y los artifacts pueden evolucionar de forma independiente, conservando mecanismos explícitos de trazabilidad. Como contrapartida, la ejecución requiere gestionar disponibilidad, permisos e integridad del almacenamiento externo.

\section*{ADR-007 --- Google Cloud como arquitectura objetivo de despliegue}
\phantomsection
\addcontentsline{toc}{section}{ADR-007 --- Google Cloud como arquitectura objetivo de despliegue}
\label{adr:007-google-cloud-despliegue}

\textbf{Estado:} aceptado; implementación DEV realizada y validada. El entorno DEMO separado y el monitoreo/alerting productivo en vivo continúan sin desplegarse.

\textbf{Contexto.} El prototipo requiere una arquitectura que permita desplegar componentes separados, almacenar artifacts de forma durable, disponer de una base PostgreSQL y controlar el consumo de recursos sin asumir inicialmente una infraestructura de alta disponibilidad propia de un entorno clínico productivo. La plataforma se planteó inicialmente como arquitectura objetivo y fue luego implementada y validada de forma progresiva.

\textbf{Decisión.} Utilizar Google Cloud como plataforma de despliegue, manteniendo servicios diferenciados para Frontend, Backend y AI Module, una instancia administrada de PostgreSQL, almacenamiento privado para modelos y assets, registro de contenedores y gestión externa de secretos. Al cierre técnico se validaron de punta a punta Artifact Registry/WIF, Backend Cloud Run DEV (legacy y Green), Frontend Cloud Run DEV, AI Cloud Run privado, Cloud SQL, Secret Manager y el storage durable, según el detalle del Capítulo 10 (Estado real del despliegue Cloud) y la evidencia técnica del Capítulo 13. El entorno DEMO separado y el monitoreo/alerting en vivo permanecen no desplegados.

\textbf{Alternativas consideradas.} Mantener el sistema exclusivamente en ejecución local, administrar infraestructura propia o trasladar la misma arquitectura a otro proveedor de nube.

\textbf{Consecuencias.} La plataforma permite utilizar servicios administrados y ajustar los recursos por componente, evitando incorporar inicialmente infraestructura que no esté justificada para el prototipo. Como contrapartida, introduce configuración específica de nube, necesidad de controlar costos y dependencia de los servicios seleccionados. El despliegue DEV validado cubre el alcance técnico definido para esta entrega; el entorno DEMO separado y el monitoreo/alerting productivo en vivo permanecen fuera del estado implementado.

\section*{ADR-008 --- Inferencia fail-closed y revision profesional obligatoria}
\addcontentsline{toc}{section}{ADR-008 --- Inferencia fail-closed y revision profesional obligatoria}

\textbf{Estado:} aceptado.

\textbf{Contexto.} Algunas capacidades de IA poseen checkpoints, metricas offline o contratos implementados, pero no cuentan con paridad runtime o preprocesamiento validado para uso automatico real.

\textbf{Decisión.} Mantener un comportamiento fail-closed: si un modelo o etapa no puede demostrar artifact, manifest, preprocesamiento y paridad runtime, el sistema debe abstenerse, bloquear o declarar estado parcial en lugar de fabricar una salida. Toda salida asistiva permanece sujeta a revision profesional obligatoria.

\textbf{Consecuencias.} Esta decision reduce el riesgo de sobreinterpretar resultados experimentales. Como contrapartida, algunas capacidades quedan visibles como bloqueadas o experimentales aunque existan componentes tecnicos parciales.

\section*{ADR-009 --- Demo Mode deterministico y no clinico}
\addcontentsline{toc}{section}{ADR-009 --- Demo Mode deterministico y no clinico}

\textbf{Estado:} aceptado.

\textbf{Contexto.} El producto necesita un recorrido reproducible para testing, defensa academica y validacion funcional, pero la localizacion anatomica real permanece bloqueada por falta de paridad runtime.

\textbf{Decisión.} Incorporar Demo Mode con fixtures deterministicos, cuatro studies y dos subjects, \texttt{source=DEMO\_FIXTURE}, \texttt{validated=false} y advertencias explicitas de no inferencia real.

\textbf{Consecuencias.} Demo Mode permite Playwright E2E y demostracion completa de producto sin afirmar inferencia clinica. No debe usarse como evidencia de validacion clinica ni como reemplazo de datos reales.

\section*{ADR-010 --- Integracion parcial del dominio Patient}
\addcontentsline{toc}{section}{ADR-010 --- Integracion parcial del dominio Patient}

\textbf{Estado:} aceptado como estado actual.

\textbf{Contexto.} Backend implementa \texttt{Patient}, \texttt{/api/patients}, \texttt{Study.patientId} y persistencia Patient--Study. El Frontend final conserva historial y longitudinal basados principalmente en \texttt{subjectRef}.

\textbf{Decisión.} Documentar Patient como dominio implementado parcialmente y evitar describirlo como gestion integral de pacientes end-to-end en browser.

\textbf{Consecuencias.} La tesis mantiene una afirmacion fiel al producto real: hay base Backend para evolucionar hacia Patient completo, pero el flujo visible actual no consume plenamente ese dominio.

\section*{ADR-011 --- No aprendizaje automatico desde correcciones profesionales}
\addcontentsline{toc}{section}{ADR-011 --- No aprendizaje automatico desde correcciones profesionales}

\textbf{Estado:} aceptado.

\textbf{Contexto.} El sistema registra revisiones, correcciones y candidatos de curacion profesional.

\textbf{Decisión.} No utilizar esas correcciones para reentrenamiento automatico ni aprendizaje online. La curacion alimenta una linea futura controlada, no un bucle activo de actualizacion de modelos.

\textbf{Consecuencias.} Se preserva reproducibilidad, trazabilidad y control de calidad de artifacts. Cualquier entrenamiento posterior requerira dataset, versionado y validacion explicitos.
